\documentclass{beamer}

% Tema
\usetheme{Madrid}
\usecolortheme{default}

% Paquetes
\usepackage[utf8]{inputenc}
\usepackage{amsmath}
\usepackage{graphicx}

% Información
\title{Módulo 3: Construcción de Escenarios Macro-Financieros}
\author{MACROPOL}
\date{\today}

\begin{document}

%------------------------------------------------
\frame{\titlepage}

%-----------------------------------------------
\begin{frame}{Contexto}
\begin{itemize}
    \item El nucleo del stress testing previsional es el escenario macro-financiero que alimenta todos los modelos.
    \item Intuicion economica:
    \begin{itemize}
        \item Los sistemas de pensiones de capitalizacion individual est\'an expuestos a:
        \begin{itemize}
            \item Ciclos econ\'omicos (crecimiento, desempleo)
            \item Condiciones fiancieras (tasas, inflaciones, retornos)
        \end{itemize}
        \item Un shcok macroeconomico se transmite a:
        \begin{itemize}
            \item Retabilidad del portafolio
            \item Densidad de cotizacoines
            \item Tasas de reemplazo
        \end{itemize}
    \end{itemize}
    \item El diseño de escenarios es el componente central del stress testing.
\end{itemize}
    
\end{frame}

%------------------------------------------------
\begin{frame}{Objetivos del módulo}
\begin{itemize}
\item Diseñar escenarios macro-financieros consistentes
\item Diferenciar baseline vs adverso
\item Modelar variables con ARIMA y VAR
\item Implementar simulación Monte Carlo
\item Traducir escenarios a inputs previsionales
\end{itemize}
\end{frame}

%------------------------------------------------
\begin{frame}{Tipos de escenarios}
\begin{tabular}{c|c}
     Tipo de escenario  &  Descripc\'on \\
     Base               & Trayectoria esperada (consistente con proyecciones oficiales \\
     Adverso            & Shock severo pero plausible
     Severo externo     & Evento                                                                                                      
\end{tabular}
    
\end{frame}
%------------------------------------------------
\begin{frame}{Arquitectura del stress testing}
\centering
Di

\begin{enumerate}
\item Escenario macro
\item Variables financieras
\item Retornos del portafolio
\item Acumulación de fondos
\item Tasa de reemplazo
\end{enumerate}
\end{frame}

%----------------------------------------
\begin{frame}{Enfoques de construccion}
\begin{itemize}
    \item Escenarios hist\'orico reproducen crisis pasadas"
    \begin{itemize}
        \item Crisis financiera global 2008
        \item  Crisis de deuda sobrerana
    \end{itemize}
    \item Tiene la ventaja del realismo
    \item Tiene la limitacion de no cubrir riesgos nuevos
    
\end{itemize}
\
    
\end{frame}


%------------------------------------------------
\begin{frame}{Tipos de escenarios}
\begin{itemize}
\item \textbf{Baseline:} trayectoria esperada
\item \textbf{Adverso:} shock plausible severo
\item \textbf{Extremo:} evento de cola
\end{itemize}
\end{frame}

%------------------------------------------------
\begin{frame}{Narrativa económica}
\textbf{Ejemplo adverso:}
\begin{itemize}
\item Crisis global
\item PIB: -3\%
\item Tasas: +300 bps
\item Bolsa: -25\%
\end{itemize}
\end{frame}

%------------------------------------------------
\begin{frame}{Variables clave}
\textbf{Macroeconómicas}
\begin{itemize}
\item PIB
\item Inflación
\item Tasa de interés
\end{itemize}

\textbf{Financieras}
\begin{itemize}
\item Acciones
\item Bonos
\item Tipo de cambio
\end{itemize}
\end{frame}

%------------------------------------------------
\begin{frame}{Enfoque metodológico}
\begin{itemize}
\item ARIMA: dinámica individual
\item VAR: interdependencia
\item Monte Carlo: distribución de resultados
\end{itemize}
\end{frame}

%================================================
\section{ARIMA}

%------------------------------------------------
\begin{frame}{Intuición ARIMA}
\begin{itemize}
\item Las series macro tienen persistencia
\item El pasado contiene información relevante
\end{itemize}
\end{frame}

%------------------------------------------------
\begin{frame}{Modelo ARIMA}
\[
y_t = c + \sum \phi_i y_{t-i} + \sum \theta_j \epsilon_{t-j} + \epsilon_t
\]
\end{frame}

%------------------------------------------------
\begin{frame}{Interpretación económica}
\begin{itemize}
\item AR: inercia
\item MA: shocks
\item Diferenciación: tendencia
\end{itemize}
\end{frame}

%------------------------------------------------
\begin{frame}[fragile]{Implementación en R}
\begin{verbatim}
library(forecast)

model <- auto.arima(inflacion)
forecast_values <- forecast(model, h=12)

plot(forecast_values)
\end{verbatim}
\end{frame}

%------------------------------------------------
\begin{frame}{Uso en pensiones}
\begin{itemize}
\item Inflación → salario real
\item Tasas → valorización de bonos
\end{itemize}
\end{frame}

%================================================
\section{VAR}

%------------------------------------------------
\begin{frame}{Intuición VAR}
\begin{itemize}
\item La economía es un sistema interconectado
\item Shock en una variable afecta a otras
\end{itemize}
\end{frame}

%------------------------------------------------
\begin{frame}{Modelo VAR}
\[
Y_t = A_1 Y_{t-1} + \dots + A_p Y_{t-p} + \epsilon_t
\]
\end{frame}

%------------------------------------------------
\begin{frame}{Interpretación}
\begin{itemize}
\item Captura interdependencias
\item Permite shocks consistentes
\end{itemize}
\end{frame}

%------------------------------------------------
\begin{frame}[fragile]{Implementación en R}
\begin{verbatim}
library(vars)

data <- cbind(pib, inflacion, tasa)
var_model <- VAR(data, p=2)

irf <- irf(var_model)
plot(irf)
\end{verbatim}
\end{frame}

%------------------------------------------------
\begin{frame}{Aplicación en pensiones}
\begin{itemize}
\item Genera trayectorias coherentes
\item Evita inconsistencias macro
\end{itemize}
\end{frame}

%================================================
\section{Monte Carlo}

%------------------------------------------------
\begin{frame}{Intuición}
\begin{itemize}
\item No hay un solo futuro
\item Se modela una distribución de resultados
\end{itemize}
\end{frame}

%------------------------------------------------
\begin{frame}{Modelo}
\[
X_{t+1} = \mu + \sigma \epsilon_t
\]
\end{frame}

%------------------------------------------------
\begin{frame}[fragile]{Implementación en R}
\begin{verbatim}
set.seed(123)

n_sim <- 1000
returns <- matrix(rnorm(n_sim*120, 
mean=0.05, sd=0.1), ncol=120)

sim_paths <- apply(returns, 1, cumsum)
\end{verbatim}
\end{frame}

%------------------------------------------------
\begin{frame}{Outputs}
\begin{itemize}
\item Percentiles (P5, P50, P95)
\item Distribución de pensiones
\end{itemize}
\end{frame}

%================================================
\section{Diseño de Escenarios}

%------------------------------------------------
\begin{frame}{Choques históricos}
\begin{itemize}
\item Crisis 2008
\item COVID-19
\end{itemize}
\end{frame}

%------------------------------------------------
\begin{frame}{Choques hipotéticos}
\begin{itemize}
\item Subida abrupta de tasas
\item Crisis externa
\end{itemize}
\end{frame}

%------------------------------------------------
\begin{frame}{Consistencia macro}
\textbf{Regla clave:}
\begin{itemize}
\item No inconsistencias (ej: PIB cae y bolsa sube fuertemente)
\end{itemize}
\end{frame}

%------------------------------------------------
\begin{frame}{Baseline vs Adverso}
\begin{tabular}{lcc}
\hline
Variable & Baseline & Adverso \\
\hline
PIB & 4\% & -2\% \\
Inflación & 3\% & 7\% \\
Bolsa & 8\% & -20\% \\
\hline
\end{tabular}
\end{frame}

%------------------------------------------------
\begin{frame}{Pipeline completo}
\begin{enumerate}
\item Estimar modelos
\item Generar shocks
\item Simular trayectorias
\item Evaluar impacto previsional
\end{enumerate}
\end{frame}

%------------------------------------------------
\begin{frame}{Ejercicio aplicado}
\textbf{Caso: República Dominicana}
\begin{itemize}
\item Estimar VAR
\item Simular shock externo
\item Generar trayectorias
\item Evaluar impacto en AFP
\end{itemize}
\end{frame}

%------------------------------------------------
\begin{frame}{Cierre}
\begin{itemize}
\item El escenario define el stress test
\item Es el vínculo entre macroeconomía y pensiones
\end{itemize}
\end{frame}

%------------------------------------------------
\end{document}